\documentclass[12pt]{ximera}
\typeout{Start loading xmPreamble.tex}%

% Add here extra macro's that are loaded automatically by all documents of claas 'ximera' or 'xourse' in this repo

%%
%%  Example:
%%
% \newcommand{\R}{\mathbb{R}}
\usepackage{tikz}
\usetikzlibrary{positioning, arrows.meta}
\usepackage{multicol}
\usepackage{enumitem}
\usepackage{caption}
\usepackage{amsmath}
\usepackage{amsthm}
\usepackage{fontenc}

%% ---------------------------------------------------------------
%% Shared worksheet macros.
%%
%% These came from the Summer 2026 worksheet set, where each worksheet carried its
%% own copy. They have to live here instead: a xourse inlines every activity into a
%% single document, so duplicate \newcommand definitions across activities collide,
%% and the xourse gobbles \input so a per-topic macro file never loads.
%%
%% They use \providecommand, NOT \newcommand. ximera.cls loads this file automatically,
%% and every activity also carries an explicit \typeout{Start loading xmPreamble.tex}%

% Add here extra macro's that are loaded automatically by all documents of claas 'ximera' or 'xourse' in this repo

%%
%%  Example:
%%
% \newcommand{\R}{\mathbb{R}}
\usepackage{tikz}
\usetikzlibrary{positioning, arrows.meta}
\usepackage{multicol}
\usepackage{enumitem}
\usepackage{caption}
\usepackage{amsmath}
\usepackage{amsthm}
\usepackage{fontenc}

%% ---------------------------------------------------------------
%% Shared worksheet macros.
%%
%% These came from the Summer 2026 worksheet set, where each worksheet carried its
%% own copy. They have to live here instead: a xourse inlines every activity into a
%% single document, so duplicate \newcommand definitions across activities collide,
%% and the xourse gobbles \input so a per-topic macro file never loads.
%%
%% They use \providecommand, NOT \newcommand. ximera.cls loads this file automatically,
%% and every activity also carries an explicit \typeout{Start loading xmPreamble.tex}%

% Add here extra macro's that are loaded automatically by all documents of claas 'ximera' or 'xourse' in this repo

%%
%%  Example:
%%
% \newcommand{\R}{\mathbb{R}}
\usepackage{tikz}
\usetikzlibrary{positioning, arrows.meta}
\usepackage{multicol}
\usepackage{enumitem}
\usepackage{caption}
\usepackage{amsmath}
\usepackage{amsthm}
\usepackage{fontenc}

%% ---------------------------------------------------------------
%% Shared worksheet macros.
%%
%% These came from the Summer 2026 worksheet set, where each worksheet carried its
%% own copy. They have to live here instead: a xourse inlines every activity into a
%% single document, so duplicate \newcommand definitions across activities collide,
%% and the xourse gobbles \input so a per-topic macro file never loads.
%%
%% They use \providecommand, NOT \newcommand. ximera.cls loads this file automatically,
%% and every activity also carries an explicit \input{../xmPreamble}, so this file is
%% read twice. That was harmless while it held only \usepackage lines (LaTeX skips a
%% package it has already loaded) but a second \newcommand is a hard error.
%%
%%   \blank[width]        fill-in-the-blank rule (default 2.5cm)
%%   \showwork            "show and explain your work" reminder
%%   \showworkline        ... plus which schedule line was used
%%   \showworkyear        ... plus which year's schedule was used
%%   \schedhead{...}      centred bold caption above a tiered-rate table
%%   \samesquares[scale]{left}{right}   two equal-area squares, labelled
%%   \samecubes[scale]{left}{right}     two equal-volume cubes, labelled
%% ---------------------------------------------------------------

\providecommand{\blank}[1][2.5cm]{\underline{\hspace{#1}}}
\providecommand{\showwork}{\textbf{REMEMBER TO SHOW AND EXPLAIN YOUR WORK.}}
\providecommand{\showworkline}{\textbf{REMEMBER TO SHOW AND EXPLAIN YOUR WORK.
  Explanation should include how and why you know which line of the
  schedule was used to calculate this amount.}}
\providecommand{\showworkyear}{\begin{sloppypar}\textbf{REMEMBER TO SHOW AND
  EXPLAIN YOUR WORK. Explanation should include how and why you know which
  line of the year's tax schedule was used to calculate this tax
  amount.}\end{sloppypar}}
\providecommand{\schedhead}[1]{\begin{center}\textbf{#1}\end{center}}

\providecommand{\samesquares}[3][1]{%
\begin{center}
{\footnotesize These squares are the \textbf{same size}.}

\vspace{1.5mm}
\begin{tikzpicture}[scale=#1,every node/.style={scale=#1}]
  \draw (0,0) rectangle (2.4,2.4);
  \node[anchor=north west] at (0.15,2.25) {Area:};
  \node[anchor=east]  at (-0.15,1.2) {#2};
  \node[anchor=north] at (1.2,-0.15) {#2};
  \begin{scope}[xshift=4.6cm]
    \draw (0,0) rectangle (2.4,2.4);
    \node[anchor=north west] at (0.15,2.25) {Area:};
    \node[anchor=east]  at (-0.15,1.2) {#3};
    \node[anchor=north] at (1.2,-0.15) {#3};
  \end{scope}
\end{tikzpicture}
\end{center}}

\providecommand{\samecubes}[3][1]{%
\begin{center}
{\footnotesize These cubes are the \textbf{same size}.}

\vspace{1.5mm}
\begin{tikzpicture}[scale=#1,every node/.style={scale=#1}]
  \foreach \dx in {0,5.2} {
    \begin{scope}[xshift=\dx cm]
      \draw (0,0) rectangle (2.0,2.0);
      \draw (0,2.0) -- (0.65,2.6) -- (2.65,2.6) -- (2.0,2.0);
      \fill[black!12] (2.0,0) -- (2.65,0.6) -- (2.65,2.6) -- (2.0,2.0) -- cycle;
      \draw (2.0,0) -- (2.65,0.6) -- (2.65,2.6) -- (2.0,2.0);
      \node[anchor=north west] at (0.1,1.9) {Volume:};
    \end{scope}
  }
  \node[anchor=east]  at (-0.15,1.0) {#2};
  \node[anchor=north] at (1.0,-0.15) {#2};
  \node[anchor=west]  at (2.25,0.4) {#2};
  \node[anchor=east]  at (5.05,1.0) {#3};
  \node[anchor=north] at (6.2,-0.15) {#3};
  \node[anchor=west]  at (7.45,0.4) {#3};
\end{tikzpicture}
\end{center}}

%% NOTE: do NOT add \usepackage to individual activity .tex files.
%% When a xourse inlines an activity it gobbles \documentclass and \input, but a bare
%% \usepackage still executes -- after \begin{document} -- and errors out.
%% All shared packages and macros belong here.
%%
%% ximera.cls already loads: amsmath, amssymb, amsthm, cancel, enumitem, graphicx,
%% hyperref, listings, pgfplots, tikz, url, xcolor. No need to re-load those.
, so this file is
%% read twice. That was harmless while it held only \usepackage lines (LaTeX skips a
%% package it has already loaded) but a second \newcommand is a hard error.
%%
%%   \blank[width]        fill-in-the-blank rule (default 2.5cm)
%%   \showwork            "show and explain your work" reminder
%%   \showworkline        ... plus which schedule line was used
%%   \showworkyear        ... plus which year's schedule was used
%%   \schedhead{...}      centred bold caption above a tiered-rate table
%%   \samesquares[scale]{left}{right}   two equal-area squares, labelled
%%   \samecubes[scale]{left}{right}     two equal-volume cubes, labelled
%% ---------------------------------------------------------------

\providecommand{\blank}[1][2.5cm]{\underline{\hspace{#1}}}
\providecommand{\showwork}{\textbf{REMEMBER TO SHOW AND EXPLAIN YOUR WORK.}}
\providecommand{\showworkline}{\textbf{REMEMBER TO SHOW AND EXPLAIN YOUR WORK.
  Explanation should include how and why you know which line of the
  schedule was used to calculate this amount.}}
\providecommand{\showworkyear}{\begin{sloppypar}\textbf{REMEMBER TO SHOW AND
  EXPLAIN YOUR WORK. Explanation should include how and why you know which
  line of the year's tax schedule was used to calculate this tax
  amount.}\end{sloppypar}}
\providecommand{\schedhead}[1]{\begin{center}\textbf{#1}\end{center}}

\providecommand{\samesquares}[3][1]{%
\begin{center}
{\footnotesize These squares are the \textbf{same size}.}

\vspace{1.5mm}
\begin{tikzpicture}[scale=#1,every node/.style={scale=#1}]
  \draw (0,0) rectangle (2.4,2.4);
  \node[anchor=north west] at (0.15,2.25) {Area:};
  \node[anchor=east]  at (-0.15,1.2) {#2};
  \node[anchor=north] at (1.2,-0.15) {#2};
  \begin{scope}[xshift=4.6cm]
    \draw (0,0) rectangle (2.4,2.4);
    \node[anchor=north west] at (0.15,2.25) {Area:};
    \node[anchor=east]  at (-0.15,1.2) {#3};
    \node[anchor=north] at (1.2,-0.15) {#3};
  \end{scope}
\end{tikzpicture}
\end{center}}

\providecommand{\samecubes}[3][1]{%
\begin{center}
{\footnotesize These cubes are the \textbf{same size}.}

\vspace{1.5mm}
\begin{tikzpicture}[scale=#1,every node/.style={scale=#1}]
  \foreach \dx in {0,5.2} {
    \begin{scope}[xshift=\dx cm]
      \draw (0,0) rectangle (2.0,2.0);
      \draw (0,2.0) -- (0.65,2.6) -- (2.65,2.6) -- (2.0,2.0);
      \fill[black!12] (2.0,0) -- (2.65,0.6) -- (2.65,2.6) -- (2.0,2.0) -- cycle;
      \draw (2.0,0) -- (2.65,0.6) -- (2.65,2.6) -- (2.0,2.0);
      \node[anchor=north west] at (0.1,1.9) {Volume:};
    \end{scope}
  }
  \node[anchor=east]  at (-0.15,1.0) {#2};
  \node[anchor=north] at (1.0,-0.15) {#2};
  \node[anchor=west]  at (2.25,0.4) {#2};
  \node[anchor=east]  at (5.05,1.0) {#3};
  \node[anchor=north] at (6.2,-0.15) {#3};
  \node[anchor=west]  at (7.45,0.4) {#3};
\end{tikzpicture}
\end{center}}

%% NOTE: do NOT add \usepackage to individual activity .tex files.
%% When a xourse inlines an activity it gobbles \documentclass and \input, but a bare
%% \usepackage still executes -- after \begin{document} -- and errors out.
%% All shared packages and macros belong here.
%%
%% ximera.cls already loads: amsmath, amssymb, amsthm, cancel, enumitem, graphicx,
%% hyperref, listings, pgfplots, tikz, url, xcolor. No need to re-load those.
, so this file is
%% read twice. That was harmless while it held only \usepackage lines (LaTeX skips a
%% package it has already loaded) but a second \newcommand is a hard error.
%%
%%   \blank[width]        fill-in-the-blank rule (default 2.5cm)
%%   \showwork            "show and explain your work" reminder
%%   \showworkline        ... plus which schedule line was used
%%   \showworkyear        ... plus which year's schedule was used
%%   \schedhead{...}      centred bold caption above a tiered-rate table
%%   \samesquares[scale]{left}{right}   two equal-area squares, labelled
%%   \samecubes[scale]{left}{right}     two equal-volume cubes, labelled
%% ---------------------------------------------------------------

\providecommand{\blank}[1][2.5cm]{\underline{\hspace{#1}}}
\providecommand{\showwork}{\textbf{REMEMBER TO SHOW AND EXPLAIN YOUR WORK.}}
\providecommand{\showworkline}{\textbf{REMEMBER TO SHOW AND EXPLAIN YOUR WORK.
  Explanation should include how and why you know which line of the
  schedule was used to calculate this amount.}}
\providecommand{\showworkyear}{\begin{sloppypar}\textbf{REMEMBER TO SHOW AND
  EXPLAIN YOUR WORK. Explanation should include how and why you know which
  line of the year's tax schedule was used to calculate this tax
  amount.}\end{sloppypar}}
\providecommand{\schedhead}[1]{\begin{center}\textbf{#1}\end{center}}

\providecommand{\samesquares}[3][1]{%
\begin{center}
{\footnotesize These squares are the \textbf{same size}.}

\vspace{1.5mm}
\begin{tikzpicture}[scale=#1,every node/.style={scale=#1}]
  \draw (0,0) rectangle (2.4,2.4);
  \node[anchor=north west] at (0.15,2.25) {Area:};
  \node[anchor=east]  at (-0.15,1.2) {#2};
  \node[anchor=north] at (1.2,-0.15) {#2};
  \begin{scope}[xshift=4.6cm]
    \draw (0,0) rectangle (2.4,2.4);
    \node[anchor=north west] at (0.15,2.25) {Area:};
    \node[anchor=east]  at (-0.15,1.2) {#3};
    \node[anchor=north] at (1.2,-0.15) {#3};
  \end{scope}
\end{tikzpicture}
\end{center}}

\providecommand{\samecubes}[3][1]{%
\begin{center}
{\footnotesize These cubes are the \textbf{same size}.}

\vspace{1.5mm}
\begin{tikzpicture}[scale=#1,every node/.style={scale=#1}]
  \foreach \dx in {0,5.2} {
    \begin{scope}[xshift=\dx cm]
      \draw (0,0) rectangle (2.0,2.0);
      \draw (0,2.0) -- (0.65,2.6) -- (2.65,2.6) -- (2.0,2.0);
      \fill[black!12] (2.0,0) -- (2.65,0.6) -- (2.65,2.6) -- (2.0,2.0) -- cycle;
      \draw (2.0,0) -- (2.65,0.6) -- (2.65,2.6) -- (2.0,2.0);
      \node[anchor=north west] at (0.1,1.9) {Volume:};
    \end{scope}
  }
  \node[anchor=east]  at (-0.15,1.0) {#2};
  \node[anchor=north] at (1.0,-0.15) {#2};
  \node[anchor=west]  at (2.25,0.4) {#2};
  \node[anchor=east]  at (5.05,1.0) {#3};
  \node[anchor=north] at (6.2,-0.15) {#3};
  \node[anchor=west]  at (7.45,0.4) {#3};
\end{tikzpicture}
\end{center}}

%% NOTE: do NOT add \usepackage to individual activity .tex files.
%% When a xourse inlines an activity it gobbles \documentclass and \input, but a bare
%% \usepackage still executes -- after \begin{document} -- and errors out.
%% All shared packages and macros belong here.
%%
%% ximera.cls already loads: amsmath, amssymb, amsthm, cancel, enumitem, graphicx,
%% hyperref, listings, pgfplots, tikz, url, xcolor. No need to re-load those.

\title{Linear Modeling: Ulna Length and Height}
\begin{document}
\begin{abstract}
The difference between a linear \emph{function} and a linear \emph{model}: taxes are
truly linear, but height-versus-ulna-length only has a linear trend. Using the class's
own measurements to build a model, then interrogating it --- what the slope and
intercept mean in context, where the domain stops making real-world sense, and how far
wrong the model is allowed to be.
\end{abstract}
\maketitle

\section*{Functions \& Models}

A \textbf{function} is an actual linear thing --- it can be used for an actual linear
scenario, such as an electric bill or a tax bracket. A \textbf{model} uses a function
to mathematically represent ``messy data.'' The data is NOT perfectly linear, but has a
linear trend which can be modeled by a linear function.

So for our ulna measurements, we are \emph{modeling} the relationship with a linear
function. With taxes, it is not a model --- the taxes are ``truly'' linear.

\section*{Modeling Ulna Data}

Before today's class, you should have measured your height and your ulna to the nearest
0.5 cm and entered that data into the spreadsheet linked in Carmen. This data was then
used to create a scatterplot and a linear model for the data.

Whether your data appears to be linear or not, record the linear function which is
modeling the data --- it will be of the form $y = m \cdot x + b$. You will refer back to
your class's own $m$ and $b$ throughout the problems below.

\begin{problem}
What does $x$ represent in the function? Use a specific value (like $x = 7$) in your
description.

\begin{freeResponse}
\end{freeResponse}

\begin{solution}
$x$ represents the length of a person's ulna, measured in centimetres. For example,
$x = 25$ refers to a person whose ulna is 25 cm long.

As with the graph-reading problems, a full answer names both the mathematical quantity
($x$, in cm) and what it means in context (the length of someone's forearm bone).
\end{solution}
\end{problem}

\begin{problem}
What does $y$ represent in the function? Use a specific value (like $y = 7$) in your
description.

\begin{freeResponse}
\end{freeResponse}

\begin{solution}
$y$ represents the predicted height of that person, in centimetres. For example,
$y = 170$ means the model predicts a height of 170 cm.

Note the word \emph{predicted}. Because this is a model rather than a function, $y$ is
what the trend line expects, not what that individual necessarily measures.
\end{solution}
\end{problem}

\begin{problem}
What is the value of the slope? What does this slope mean in terms of this context?

\begin{freeResponse}
\end{freeResponse}

\begin{solution}
The slope is the $m$ from your class's model. In context it says: for every additional
1 cm of ulna length, the model predicts roughly $m$ additional centimetres of height.

Published studies put this figure somewhere near 4, so if your class's slope is in that
neighbourhood, an extra centimetre of forearm corresponds to about four extra
centimetres of person.
\end{solution}
\end{problem}

\begin{problem}
What is the $y$-intercept value? What does this mean in terms of this context?

\begin{freeResponse}
\end{freeResponse}

\begin{solution}
The $y$-intercept is the $b$ from your class's model. Literally, it is the height the
model predicts for a person whose ulna is 0 cm long --- which is why it is worth
examining rather than accepting.
\end{solution}
\end{problem}

\subsection*{Establishing a Reasonable Domain}

\begin{problem}
Does the $y$-intercept's contextual meaning ``make sense'' in the real world?
\begin{multipleChoice}
\choice{Yes}
\choice[correct]{No}
\end{multipleChoice}

\begin{freeResponse}
\end{freeResponse}

\begin{solution}
No. The $y$-intercept describes a person with an ulna of length 0 cm, which does not
describe any person at all. The model still needs an intercept to position the line
correctly, but that particular point carries no real-world meaning.

This is the central lesson about models: the algebra will happily give you an output
for any input you feed it, and it is your job to decide which inputs are meaningful.
\end{solution}
\end{problem}

\begin{problem}
Would the meaning for $x = 5$ ``make sense'' in the real world?
\begin{multipleChoice}
\choice{Yes}
\choice[correct]{No}
\end{multipleChoice}

Would the meaning for $x = 50$ ``make sense'' in the real world?
\begin{multipleChoice}
\choice{Yes}
\choice[correct]{No}
\end{multipleChoice}

\begin{solution}
Neither does. A 5 cm ulna is far shorter than even an infant's, and a 50 cm ulna is far
longer than any adult's --- adult ulnae run roughly 24--30 cm. Both inputs sit well
outside the range of the data the model was built from, so the model has no grounds to
say anything about them.
\end{solution}
\end{problem}

\begin{problem}
What do you think would be the smallest $x$ value for which our model would give a
reasonable real-world interpretation? Explain your choice in 2--3 sentences.

\begin{freeResponse}
\end{freeResponse}

\begin{solution}
A defensible answer sits near the smallest ulna measurement your class actually
recorded. The reasoning matters more than the number: the model was fitted to your
class's data, so it is only supported over the range that data covers. Going below the
smallest measured value is extrapolation, and the further you extrapolate the less
reason you have to trust the output.
\end{solution}
\end{problem}

\begin{problem}
What do you think would be the largest $x$ value for which our model would give a
reasonable real-world interpretation? Explain your choice in 2--3 sentences.

\begin{freeResponse}
\end{freeResponse}

\begin{solution}
Symmetrically, a defensible answer sits near the largest ulna measurement in your
class's data, for the same reason.

The two values you have now chosen --- the smallest and largest sensible $x$ --- are
together the \textbf{domain} of the model.
\end{solution}
\end{problem}

\section*{The Model's Error}

\begin{problem}
Suppose someone's ulna bone is 24 cm long. According to our model, how tall is this
person likely to be?

\begin{freeResponse}
\end{freeResponse}

\begin{solution}
Substitute $x = 24$ into your class's model: the predicted height is $m \cdot 24 + b$
centimetres. Using the researchers' published model from the next section,
$y = 4.39(24) + 45.89 = 151.25$ cm, so a class model in that neighbourhood should give
something close to 151 cm.
\end{solution}
\end{problem}

\begin{problem}
Was there a data point with $x = 24$? If so, how does it compare to the output you got
in the previous question? If there is not a data point with $x = 24$, estimate how far
off you think the model could be from a ``true'' value with $x = 24$, based on how the
model relates to other data points.

\begin{freeResponse}
\end{freeResponse}

\begin{solution}
This depends on your class's data. The thing to look at is the vertical gap between the
data points and the line near $x = 24$: if the nearby points sit within a few
centimetres of the line, that is roughly how far off you should expect the prediction at
$x = 24$ to be.
\end{solution}
\end{problem}

\begin{problem}
What do you think is an acceptable amount that the model's height could be off? Is off
by 5 cm acceptable? Off by 10 cm? By 1 meter? Explain your choice in 2--3 sentences.

\begin{freeResponse}
\end{freeResponse}

\begin{solution}
There is no single right answer, but the reasoning should be anchored to the context. 1
metre is plainly unacceptable --- that is the difference between a child and a tall
adult. 5 cm is about two inches, which for estimating someone's height from a forearm
bone is arguably fine. 10 cm is around four inches, which starts to be a meaningful
error.

What makes an answer good here is naming what the model would be \emph{used for}. A
forensic identification demands far tighter error than a rough classroom estimate.
\end{solution}
\end{problem}

\begin{problem}
Suppose someone is 170 cm tall. About how long is their ulna bone likely to be
according to your model?

\begin{freeResponse}
\end{freeResponse}

\begin{solution}
This time you know $y$ and want $x$, so solve $170 = m \cdot x + b$ for $x$, giving
$x = \dfrac{170 - b}{m}$.

With the researchers' model, $x = \dfrac{170 - 45.89}{4.39} = \dfrac{124.11}{4.39}
\approx 28.3$ cm.
\end{solution}
\end{problem}

\begin{problem}
Will this ulna value be close to the ``true'' ulna length of a person who is 170 cm
tall? Why or why not?

\begin{freeResponse}
\end{freeResponse}

\begin{solution}
Not necessarily. Plenty of different people are 170 cm tall, and their ulnae are not all
the same length --- that is exactly why the data forms a cloud rather than a perfect
line. The model returns the ulna length of a \emph{typical} 170 cm person, not of any
particular one.
\end{solution}
\end{problem}

\begin{problem}
Will its ``error'' be the same as the ``acceptable error'' you established earlier?

\begin{freeResponse}
\end{freeResponse}

\begin{solution}
No, and this is the subtle point. Earlier you set an acceptable error in \emph{height},
measured in centimetres of height. Now you are predicting an \emph{ulna length}, so the
error is measured in centimetres of ulna.

Because the slope is about 4, an error of 4 cm in predicted height corresponds to only
about 1 cm of ulna. The two errors are related by the slope, so a tolerance you set for
one does not transfer directly to the other.
\end{solution}
\end{problem}

\section*{A Published Scientific Model}

Below is a graph of data from a study which measured the ulna and height of many women
to see if there is a linear relationship between the two.

\begin{center}
\includegraphics[width=\textwidth]{linearModels/ulna_scatterplot.png}
\end{center}

The researchers decided that this data does have a linear relationship, and found the
linear trendline for this data to be
\[
y = 4.39x + 45.89.
\]

\begin{problem}
Describe how you would draw this line on the graph above. Note that you are not
``sketching'' the line; you have an expression for the function, so you should be able
to place it fairly exactly.

(Hint: on a separate sheet of paper, draw a line which passes through $(3,4)$ and
$(6,10)$. Can you do something similar here with the information provided?)

\begin{freeResponse}
\end{freeResponse}

\begin{solution}
Pick two $x$-values inside the range of the data and compute the matching $y$-values,
then join them. For instance at $x = 24$, $y = 4.39(24) + 45.89 = 151.25$; at $x = 30$,
$y = 4.39(30) + 45.89 = 177.59$. Plot $(24, 151.25)$ and $(30, 177.59)$ and draw the
line through them.

Two points determine a line exactly, which is why this is drawing rather than sketching.
\end{solution}
\end{problem}

\begin{problem}
Along with this trendline data, the researchers specified that the resulting $y$-value
from this trendline could be $\pm 7.03$ in value. What does this mean in terms of this
context?

\begin{freeResponse}
\end{freeResponse}

\begin{solution}
It means the model's predicted height for a given ulna length is expected to be within
about 7.03 cm of that person's actual height. So for an ulna of 24 cm, predicting
151.25 cm, the person's real height would typically fall somewhere between about 144.2
cm and 158.3 cm.

That is a band roughly 14 cm wide --- worth noticing when you compare it to the
acceptable error you chose earlier.
\end{solution}
\end{problem}

\begin{problem}
Describe how you would draw two more lines on the graph to represent this $\pm 7.03$
information, and explain how you used the information to decide where they go.

\begin{freeResponse}
\end{freeResponse}

\begin{solution}
Draw one line 7.03 units above the trendline and one 7.03 units below it, both parallel
to the original. In equation form these are $y = 4.39x + 52.92$ and
$y = 4.39x + 38.86$ --- the slope is unchanged because the uncertainty is in the
predicted height, not in the rate of change, so the band has constant vertical width.

Most of the data points should fall inside this band.
\end{solution}
\end{problem}

\begin{problem}
Pick a data point that might describe an individual who has ``proportionally short
arms'' and describe where it sits on the graph. What are you looking for graphically,
and why?

\begin{freeResponse}
\end{freeResponse}

\begin{solution}
Look for a point sitting well \emph{above} the trendline. A point above the line is
someone whose actual height is greater than the model predicts from their ulna ---
equivalently, their ulna is shorter than you would expect for someone of their height.
That is what ``proportionally short arms'' means.
\end{solution}
\end{problem}

\begin{problem}
Pick a data point that might describe an individual who has ``proportionally long arms''
and describe where it sits on the graph. What are you looking for graphically, and why?

\begin{freeResponse}
\end{freeResponse}

\begin{solution}
Look for a point sitting well \emph{below} the trendline: someone whose actual height is
less than the model predicts from their ulna, meaning their forearm is long relative to
their height.

The vertical distance from a point to the line is called the \emph{residual}, and it is
precisely this quantity --- how far off the model is for that individual --- that the
$\pm 7.03$ figure summarises across the whole dataset.
\end{solution}
\end{problem}

\end{document}
